\subsection{Places}

If K is a field, recall that \(\tilde{K}\) denotes the set the sum of K and an element denoted by \(\infty\) (Algebra, Chapter II, § 9, no. 9); the laws of composition of K extend to \(\tilde{K}\) by setting (loc. cit.)

\begin{enumerate}
\def\labelenumi{(\arabic{enumi})}
\setcounter{enumi}{1}
\tightlist
\item
\end{enumerate}

\[
a + \infty = \infty \quad \text { for } \quad a \in \mathbf {K}, \quad a \neq \infty ,\tag{2}
\]

\[
\infty . a = a. \infty = \infty \quad \text {   for   } \quad a \in \tilde {\mathbf {K}}, \quad a \neq 0.\tag{3}
\]

The only compositions not defined are therefore the compositions \(\infty + \infty\) , \(\infty, 0\) and \(0, \infty\) . On the other hand, the mappings \(x \mapsto -x\) and \(x \mapsto x^{-1}\) extend similarly to \(\tilde{\mathbf{K}}\) by setting \(-\infty = \infty, 0^{-1} = \infty, \infty^{-1} = 0\) . We shall also write \(x + (-y) = x - y\) .

The set \(\tilde{\mathbf{K}}\) , called the projective field associated with \(\mathbf{K}\) , can be identified with the projective line \(\mathbf{P}_1(\mathbf{K})\) (loc. cit.).

DEFINITION 3. Let K and L be two fields. Every morphism f of \(\tilde{\mathbf{K}}\) to \(\tilde{\mathbf{L}}\) (for addition and multiplication) such that \(f(1) = 1\) is called a place of K with values in L.

In other words, if \(x\) and \(y\) are elements of \(\tilde{\mathbf{K}}\) and \(f(x) + f(y)\) (resp. \(f(x)f(y)\) ) is defined, then \(x + y\) (resp. \(xy\) ) is defined and

\begin{enumerate}
\def\labelenumi{(\arabic{enumi})}
\setcounter{enumi}{3}
\tightlist
\item
\end{enumerate}

\[
f (x + y) = f (x) + f (y)\tag{4}
\]

\[
f (x y) = f (x) f (y).\tag{5}
\]

As \(\infty +\infty\) is not defined, neither is \(f(\infty) + f(\infty)\) , which shows that

\[
f (\infty) = \infty .\tag{6}
\]

Similarly, since \(0.\infty\) is not defined, neither is \(f(0)f(\infty)\) , which, by virtue of (6), implies

\[
f (0) = 0.\tag{7}
\]

On the other hand

\[
f (a ^ {- 1}) = f (a) ^ {- 1} \quad \text {   for   all   } a \in \tilde {\mathbf {K}}.\tag{8}
\]

If \(f(a)f(a^{-1})\) is defined, \(aa^{-1}\) is defined and hence is equal to 1; then \(f(a)f(a^{-1}) = f(1) = 1\) , which proves (8) in this case. If \(f(a)f(a^{-1})\) is not defined, then, either \(f(a) = 0\) and \(f(a^{-1}) = \infty\) or \(f(a) = \infty\) and \(f(a^{-1}) = 0\) and (8) still holds.

Similarly it can be shown that

\[
f (- a) = - f (a) \quad \text {   for   all   } a \in \tilde {\mathbf {K}}.\tag{9}
\]

From formulae (8) and (9) it follows that \(f\) is also a morphism for the laws of composition \((x, y) \mapsto x - y\) and \((x, y) \mapsto xy^{-1}\) .

For \(x \in \tilde{\mathbf{K}}\) , \(f\) is called finite at \(x\) if \(f(x) \neq \infty\) ; this implies \(x \in \mathbf{K}\) by (6).

If \(f \colon \tilde{\mathbf{K}} \to \tilde{\mathbf{L}}\) is a place, \(f(\tilde{\mathbf{K}})\) is a subset of \(\tilde{\mathbf{L}}\) which is stable for the laws of composition \((x, y) \mapsto x + y\) , \((x, y) \mapsto x - y\) , \((x, y) \mapsto xy\) and \((x, y) \mapsto xy^{-1}\) and which contains 1. If \(\mathbf{E}\) is the set of finite elements of \(f(\tilde{\mathbf{K}})\) , \(\mathbf{E}\) is a subfield of \(\mathbf{L}\) and \(f(\tilde{\mathbf{K}}) = \tilde{\mathbf{E}}\) . By an abuse of language \(\mathbf{E}\) is called the value field of \(f\) .

The composite mapping of two places is a place.

Let F be an isomorphism of a field K onto a subfield of a field L; let us extend f to \(\tilde{K}\) by setting \(f(\infty) = \infty\) . Thus we obtain a place of K with values in L which is called trivial and which is often identified with the isomorphism f.

